% Chapter 8 — Ingestion, Corporate Actions, and the Market Data Operator (NAZAROV, 8.5 pp).
% Discharges Constitution v1.1 §9. Carries G1 (four mechanisms + spin-off/elective-merger
% half-page examples), G3 (W1-W4), G5 (cum/ex). Declares the boundary-convention slot Ch.12 consumes.
\section{Ingestion, Corporate Actions, and the Market Data Operator}\label{ch:marketdata}

This chapter discharges \S 9 of the constitution.

\S 9 governs the one place the closed system is not closed: the boundary where external
data crosses in. Everything inside the ledger conserves quantity by construction
(Chapter~\ref{ch:objects}); value, by contrast, is supplied from outside --- prices,
dividend figures, index levels, fixings, corporate-action notices, stamped valuations.
The conservation properties hold only because the system is closed, so the boundary is
where they are won or lost. A datum admitted without provenance, a corporate action that
silently re-reads old data in a new frame, a fallback taken without a record: each
corrupts the closed-system property without violating any single internal invariant. This
chapter holds the boundary watertight, and it does so with two devices --- the recorded
observation and the market data operator --- and no third.

\subsection{The boundary, and corporate actions on it}

Two things cross the boundary, and they are recorded the same way. \emph{Market and
reference data} --- a price, a dividend figure, an index level, a fixing, a stamped NAV
--- crosses as a recorded observation. A \emph{corporate action} --- a split, a dividend,
a spin-off, a merger --- crosses first as a recorded notice and is then, on its effective
date, a transaction like any other: a bundle of moves, unit-state changes, and a new
version of the instrument's terms, applied atomically through the one Transaction Executor
(Chapter~\ref{ch:machines}). A corporate action is not a privileged operation on a store;
it is admitted at the same door as a trade, and it writes ProductTerms, UnitStatus, and
balances through the same four-kind transaction (Chapter~\ref{ch:objects}).

A corporate action does one thing no trade does: it changes the meaning of every datum
recorded before it. After ACME's two-for-one split, the recorded close $100.00$ still
names the same economic fact, but a post-split reader who multiplies it by the post-split
count $20{,}000$ books a phantom. To carry a pre-event datum into the post-event frame a
transformation must be applied, and the framework gives that transformation a name: the
\textbf{market data operator}. Three principles govern it, taken from \S 9 and used
throughout this chapter:

\begin{enumerate}[label=(\roman*),nosep,leftmargin=2.2em]
\item the operator is \emph{intrinsic to the pairing} of a kind of data with a kind of
  event, declared once when the kind of data is introduced and the corporate action is
  recorded, never improvised at the moment of reading;
\item the \emph{ledger is the authority} and the pricing data layer merely applies: the
  operator is declared data on the record, and the pricing data layer computes with it,
  never around it;
\item \emph{originals are never overwritten}. The observed datum stays on the record
  exactly as observed; the adjusted value is computed at each read; derived quantities are
  recomputed from operator-adjusted inputs --- until fresh post-event data are recorded,
  after which the operator has nothing left to adjust.
\end{enumerate}

\subsection{The observation door}\label{sec:obs-door}

No datum enters the ledger as a bare read. Every datum crosses as a recorded
\textbf{observation} carrying its provenance: the source that published it, the time it
was observed, and the basis frame it was observed in. This is the door that
Chapters~\ref{ch:valuation} and~\ref{ch:virtual} refer to when a mark or a stamped level
``enters through the observation door.'' A read against a live feed leaves no record and
is therefore unreproducible; a recorded observation is a fact of the log, so every
downstream number is a function of the record alone (\S 12).

\begin{lstlisting}
data Provenance  = Provenance { source :: SourceId, observed :: Time, basis :: BasisRef }
data Observation = Observation { datum :: Datum, prov :: Provenance }
-- every datum crosses as a recorded Observation; there is no bare read.
-- the door admits or refuses; a refusal is a returned value, never a partial write:
data Ingest      = Admitted Observation | Refused IngestError
\end{lstlisting}

The provenance frame is what makes ``prices and quantities are never mixed across an
event'' (\S 8) enforceable rather than hoped for: the frame a datum was observed in is on
the record beside the datum, so a read that would combine it with a quantity in a
different frame is detectable, not silent. That detection is the invariance witness of
Section~\ref{sec:witness}.

\medskip\noindent\textbf{Episode --- the variance swap (T5): the fixing source is the
record.}
\emph{Trigger.} VS-1's scheduled fixing dates arrive; each is a watch on the recorded IDX
official close (Chapter~\ref{ch:objects}). \emph{The door.} The fixing is not a second
feed; it is a read of the same recorded official-close observation that FUT-IDX settles
against. The three settlement days of the future --- IDX closes $1{,}000.00$, $990.00$,
$1{,}005.00$ --- are fixings $124$, $125$, and $126$ of the swap's $252$-fixing series.
\emph{Design point.} \emph{One recorded observation source feeds two consumers; because
the source is on the record, the fixing series is deterministic and replayable at the data
boundary, and firing $127$ recomputes the same statistic any party would
(Chapter~\ref{ch:valuation}).}

\medskip\noindent\textbf{Episode --- the total return swap (T6): the bridge is an
observation.}
\emph{Trigger.} TRS-1's first quarterly reset arrives; virtual ledger V-1's replayable
level is stamped at \USD{10{,}300{,}000}. \emph{The door.} That level crosses as an
observation with provenance --- source the fold of V-1's own log, observed at the reset
date, in the reference-currency frame --- and is recorded. No move crosses: a move from a
wallet of V-1 to a wallet of the true ledger has no paired leg that conserves in either
ledger, so neither Transaction Executor can admit it (Chapter~\ref{ch:virtual}). \emph{The
door checks.} Provenance present; the crossing is a recorded observation, not a move.
\emph{Design point.} \emph{The bridge between the virtual ledger and the true ledger is an
observation, never a move; the reset and financing moves it triggers are the true ledger's
own (Chapter~\ref{ch:virtual}).}

\subsection{The datum-kind registry}\label{sec:registry}

Market data is not one number per name. A dividend forecast is a composite: cash entries
that carry a price dimension and scale under a split, and proportional entries that carry
no dimension and stand. An operator that treated the datum as a single scalar would either
scale the proportional part, which is wrong, or leave the cash part, which is also wrong.
The operator must therefore know the \emph{dimension of each component}, and it learns it
from the \textbf{datum-kind registry}.

A kind is registered once, before any datum of that kind crosses, with a component-wise
declaration: each component's dimension --- price-in-basis, proportional, or quantity ---
and the quotation convention it renders. Registration is immutable: a correction is a new
kind plus a retirement of the old, never an edit, so no datum ever changes dimension under
a reader's feet. A datum whose kind is not registered is refused at the door.

\begin{lstlisting}
data Dim      = Price | Proportional | Quantity          -- component dimensions
data Datum    = Datum { kind :: KindId, comps :: [(Dim, Integer)] }  -- exact minor units
data IngestError = UnknownUnit | UndeclaredConvention | UnregisteredKind
-- an operator acts on Price components; Proportional stands; Quantity is walled off
-- to the move plane (Chapter 3) and is never carried by an operator.
\end{lstlisting}

The wall between \texttt{Quantity} and the operators is the same wall as the move plane. A
share count is not adjusted by an operator: the two-for-one split moves $10{,}000$ ACME to
$20{,}000$ by a paired-leg move from the issuer's virtual wallet (Chapter~\ref{ch:objects}),
because quantity is conserved and only a move may change it. Operators act on value-frame
components only. This division is what makes the split's arithmetic come out invariant, and
it is stated as the witness next.

The registry is a named trust assumption, and it is the only one the operators rest on.

\begin{quote}
\textbf{TA-KIND.} A kind declaration faithfully renders the publisher's quotation
convention --- which components are cash and which are proportional, in which basis.
\emph{Owner:} the party that governs kind registration (data governance). \emph{Violation:}
an operator scales a component it should not, or spares one it should scale, misvaluing
across every subsequent corporate action. \emph{Detection:} the invariance witness at the
first crossing of a corporate action (Section~\ref{sec:witness}), divergence against an
independent source (regime W3 below), and \S 13 recomputation. \emph{Residual:} a
misdeclaration that never meets a corporate action is invisible from inside the boundary
--- a reconciliation matter at the perimeter, the exact sibling of the collateral regime
bit of Chapter~\ref{ch:collateral}.
\end{quote}

\subsection{The market data operator and adjustment-schedule totality}\label{sec:operator}

An operator is a function of the registered kind and the event kind, declared once and
recorded. It carries a value-frame datum from the pre-event basis to the post-event basis;
it is the identity on proportional components and undefined on quantity components, which
never reach it.

\begin{lstlisting}
-- one operator per (event kind, datum kind), declared once, never improvised at read:
operator :: EventKind -> KindId -> (BasisFrame -> BasisFrame)
\end{lstlisting}

The first mechanism is that this declaration is \emph{total}.

\begin{principle}[Adjustment-schedule totality]\label{prin:sched-total}
For every registered datum kind and every event kind the boundary admits, an operator is
declared --- the identity operator where a kind is genuinely unaffected, but declared, not
assumed. There is no (datum kind, event kind) pair for which a read finds the adjustment
undefined. A corporate action whose pairing has no declared operator does not cross: it is
refused at the door, never silently treated as identity.
\end{principle}

Totality is what turns ``never improvised at read'' into a checkable property. An
improvised adjustment is precisely a pairing with no declared operator; forbidding the
improvisation means requiring the declaration in advance, for the whole registered domain.
The refusal of an undeclared pairing is failure regime W4 (Section~\ref{sec:failures}), and
it fails closed, as \S 2's Correctness commitment demands: the system stops rather than
guesses a frame.

Operators compose. Two corporate actions in sequence adjust a datum by the composition of
their operators, and the composition is taken \emph{in log order}, because the operators do
not commute: a split then a special dividend is not the same adjustment as the dividend
then the split. The log fixes the order, so the composed adjustment is well defined and
replayable.

\begin{secondtelling}
Fix a datum kind. Its bases are the objects of a small category; each recorded event
contributes the arrow its operator names, from the pre-event basis to the post-event basis;
identity events contribute identity arrows. A chain of corporate actions is then a path, and
the adjustment it induces is the composite arrow --- associative because composition of
functions is, order-sensitive because the arrows do not commute. The map sending each event
sequence to its composed operator is a functor from event paths to basis maps; adjustment
totality is the statement that this functor is defined on every arrow of the source. Nothing
below depends on this telling.
\end{secondtelling}

\subsection{The operator menu: the split}\label{sec:menu}

The split shows the whole menu at once. On 2026-06-01 ACME's two-for-one split is
effective; the recorded split notice had declared the effective-date watch, the Event
Monitor emits the event, and one transaction carries the move and the new terms version
(Chapters~\ref{ch:objects} and~\ref{ch:homes}). What each recorded datum kind reads as, in
the new frame, is the operator's work:

\begin{center}
\begin{tabular}{@{}lll@{}}
\toprule
Datum kind & Operator for the 2-for-1 split & Read after the event \\
\midrule
Official close (spot) & scale by $\tfrac12$ & $100.00 \rightarrow 50.00$ \\
Cash-per-share dividend figure & scale by $\tfrac12$ & $1.50 \rightarrow 0.75$ \\
Proportional (percentage) figure & identity --- stands & $1.00\% \rightarrow 1.00\%$ \\
Index divisor (a constituent splits) & adjusted to hold the level & level continuous \\
\bottomrule
\end{tabular}
\end{center}

\noindent The share count is absent from the menu on purpose: $10{,}000 \rightarrow
20{,}000$ is a move, not an operator (Section~\ref{sec:registry}). The spot and the
cash-per-share figure are price-dimensioned and halve; the percentage figure is
proportional and stands; the index divisor is adjusted so the published level does not jump
when a constituent splits. Every original stays on the record; each ``read after'' is
computed at the read, until a fresh post-split close prints and the operator retires for
that datum.

\medskip\noindent\textbf{Episode --- the dividend figure under the split (T3$\times$T4).}
\emph{Trigger.} On 2026-05-04 ACME's cash dividend of \USD{1.50} per share is announced;
the announcement is captured at the boundary as a recorded observation, and its declared
per-share figure $1.50$ enters the record (the three firings that pay it are
Chapter~\ref{ch:contracts}'s). \emph{The operator.} The split is effective 2026-06-01,
after the payment date but within the window in which the declared forward figure is still
read; a read of the $1.50$ figure in the post-split frame applies the split's
cash-per-share operator and returns $0.75$. \emph{The door checks.} The datum kind is
registered as a cash-per-share (price) component, so the operator scales it; a proportional
figure beside it would have stood. \emph{Design point.} \emph{One operator, one invariant,
visible: the same halving that takes the spot from $100.00$ to $50.00$ takes the recorded
per-share dividend figure from $1.50$ to $0.75$, so the dividend history reads in
current-share terms --- the scaling of a real figure is shown, not asserted.}

\subsection{State-aware reads: consistency of reference and cum/ex}\label{sec:witness}

Every operator read is state-aware: the frame a datum reads in is determined by the
corporate-action state of its underlying at the read, and that state is on the record. Two
readings of one recorded datum in two frames are two different values, and combining a
quantity with a price drawn from a different frame produces a phantom. The check that
refuses phantoms is the \textbf{invariance witness}: the operator on the value frame and
the move on the quantity are bound so their product --- total value --- is invariant across
the event.

For the split the witness is arithmetic. Before: $10{,}000 \times 100.00 = \USD{1{,}000{,}000}$.
After, read in the post-split frame: $20{,}000 \times 50.00 = \USD{1{,}000{,}000}$. The
phantom the witness refuses is $20{,}000 \times 100.00 = \USD{2{,}000{,}000}$, the new count
against the stale frame. This is the consistency-of-reference invariant; it is stated
formally and catalogued in Chapter~\ref{ch:invariants}, and used here.

The same witness governs the cum/ex boundary of a distribution, where no quantity changes
at all. Before the ex date a price is \emph{cum} --- it includes the coming distribution;
on and after the ex date it is \emph{ex} --- it excludes it. The cum price and the ex price
are one datum kind read in two states separated by the ex-date event, whose declared
operator subtracts the distribution's value. A cum price consumed as though it were ex, or
an ex price consumed as though it were cum, is a phantom of exactly the split's kind ---
the right number in the wrong frame --- and the same invariance witness refuses it, because
the state that fixes the frame is the recorded ex-date event, not the wall clock. A read
that ignores the ex-date state prices the instrument on the wrong side of its own corporate
action, and the witness catches it before the number reaches a valuation.

\subsection{Recompose and the aggregate weld}\label{sec:recompose}

The split adjusts one unit's data in place. Two further mechanisms handle corporate actions
that change which units exist: one splits a unit into several, the other merges several into
one. Both preserve total value, and both are checked by the same witness.

\paragraph{Recompose (one unit becomes many): a spin-off.}
When an issuer distributes a new entity to its holders, one position becomes a basket, and
the pre-event data of the parent must be carried onto the resulting units. \emph{Recompose}
is the operator that does so, per a declared allocation basis recorded with the event.
Illustratively: HELIOS trades at $120.00$ and spins off LUMEN, one LUMEN share per HELIOS
share, with a declared allocation of $75\%$ of pre-event value to the HELIOS stub and
$25\%$ to LUMEN. A holder of $1{,}000$ HELIOS holds, after the event, $1{,}000$ HELIOS and
$1{,}000$ LUMEN --- the new LUMEN mass arriving by a paired-leg move from the issuer, not by
an operator. Recompose carries the recorded pre-event close onto the two frames: HELIOS stub
$90.00$, LUMEN $30.00$. The witness holds: $1{,}000 \times 120.00 = 1{,}000 \times 90.00 +
1{,}000 \times 30.00 = \USD{120{,}000}$. The original $120.00$ is preserved; the stub and
spun-off frames are computed at read until fresh post-event closes print for each.

\paragraph{The aggregate weld (many units become one): an elective merger.}
When one unit is absorbed into another, the resulting position welds into the acquirer's
aggregate, and the two data histories must weld into one continuous series so that valuation
and profit and loss run unbroken across the event. The \emph{aggregate weld} is the operator
that carries the pre-event data of the absorbed unit onto the acquirer's frame, per the
elected terms recorded at the boundary. An elective merger is captured as an observation ---
the election result is a datum with provenance, like any other. Illustratively: TARGET is
acquired by ACQUIRER, each TARGET share electing either \USD{50.00} cash or $0.4$ ACQUIRER
shares; ACQUIRER trades at $125.00$. A holder of $1{,}000$ TARGET who elects stock welds to
$400$ ACQUIRER, and the pre-merger TARGET series is welded onto the ACQUIRER frame at the
elected ratio $0.4$; the witness holds, $1{,}000 \times 50.00 = 400 \times 125.00 =
\USD{50{,}000}$. A holder who elects cash has the position retired at settlement and no weld
occurs --- the elected terms, recorded data, decide which mechanism runs.

Two further composite distributions --- a dividend forecast and a special dividend --- are
governed by the same registry and the same operators: the forecast is a composite of cash
and proportional components that the split scales component-wise, and the special dividend
carries a cum/ex operator like the ordinary one. Their fully worked arithmetic is deferred
to the Worked-Examples companion (register lines~E71 and~E72); the mechanism here is
complete without them.

\subsection{Failure regimes at the boundary}\label{sec:failures}

The boundary has four failure modes, and each has a determined, replayable disposition --- a
function of log state alone, so replay reproduces it bit-identically. None improvises; each
fails closed or records a flag, never a silent pick.

\begin{center}
\begin{tabular}{@{}lll@{}}
\toprule
 & Boundary condition & Determined, replayable disposition \\
\midrule
W1 & missing datum & the operator applies to the last recorded original, carrying a \\
   &               & staleness flag; with no original, the price is undefined, not \\
   &               & zero, and valuation over the unit defers (Chapter~\ref{ch:valuation}) \\
W2 & late datum & recorded at its observation-time coordinate; values derived over \\
   &            & $[\text{observed},\text{arrival})$ are recomputed --- ``as known at $t$'' and \\
   &            & ``with corrections through $t'$'' are both queryable \\
W3 & contradictory sources & sources diverging beyond the declared threshold yield a \\
   &                       & flagged ``aggregation failed'' observation, never a silent \\
   &                       & pick; sources and threshold are recorded \\
W4 & unmappable event & a pairing with no declared operator, or a datum of an \\
   &                  & unregistered kind, is refused at the door (Principle~\ref{prin:sched-total}) \\
\bottomrule
\end{tabular}
\end{center}

\noindent W1 is the working state of \S 9's ``until fresh post-event data are recorded'':
between the event and the first fresh datum, every read is the original under its operator,
flagged stale, and the flag is a fact of the record, not a property of a cache. W2 makes the
boundary bitemporal: because an observation carries its own observation time, a datum
arriving late does not rewrite history --- it is placed at its own time, and the derived
folds over the affected interval are recomputed, both the original and the corrected view
surviving (\S 2, \S 12). W3 is the multi-source case: a single source is a single point of
failure, so where a datum materially affects valuation it is drawn from more than one
attested source, aggregated by a declared rule, and a divergence beyond the declared
threshold produces a flagged output rather than a silently chosen one. W4 is
adjustment-schedule totality enforced: the undeclared pairing cannot cross.

\subsection{The boundary-convention slot}\label{sec:slot}

One class of boundary fact is neither a datum nor an operator: a \emph{convention} that
some downstream projection needs and that must be stated once, as data, so that two readers
of one log compute one answer. \S 9 places such conventions on the record beside the
operators; this section declares the \textbf{slot} they occupy, and
Chapter~\ref{ch:reporting} fills and consumes one.

A convention in the slot is a total function of log state. The convention
Chapter~\ref{ch:reporting} consumes --- the attribution of commingled received value across
a reinvestment pool --- has three components, and the slot's shape is exactly these three:

\begin{lstlisting}
data Convention = Convention
  { pool     :: Predicate (Wallet, Move)   -- the pool over which value is attributed
  , basis    :: PeriodBasis                -- when the pro-rata ratio is struck
  , rounding :: RoundingRule }             -- remainders booked explicitly (Chapter 3)
attribute :: Convention -> LogPrefix -> Partition Amount   -- consumed in Chapter 12
\end{lstlisting}

The slot declaration is itself a recorded event, so it versions with the log prefix: a
report computed as of a past date uses the convention then in force, and a changed
convention is a new recorded version, never an edit. Totality is the requirement that makes
the slot safe to consume: for every log prefix the convention assigns every attributable
quantity, the attributed parts and the explicit remainder partition the whole, and the
rounding rule leaves no quantity silently dropped (\S 4's remainder discipline). A
convention missing any component is refused, because two readers would then compute
different figures from one log --- internal reconciliation failure reintroduced through a
report. Chapter~\ref{ch:reporting} states the totality property and proves the figures it
builds are projections of coordinates, units, and this declared convention alone; this
chapter's duty is discharged by declaring the slot and its totality obligation.

\subsection{What an auditor checks}

A candidate implementation satisfies this chapter when six things hold, each mechanically
checkable over generated data and histories: every datum on the record carries provenance,
and no valuation reads a datum that does not (Section~\ref{sec:obs-door}); every datum's
kind is registered with component dimensions before it crosses, and an unregistered kind is
refused (Section~\ref{sec:registry}); an operator is declared for every (registered kind,
event kind) pair, and an undeclared pairing is refused, not defaulted to identity
(Principle~\ref{prin:sched-total}); every original is preserved and every adjusted value is
recomputed at read from the original under the declared operators
(Sections~\ref{sec:menu}--\ref{sec:recompose}); the invariance witness holds across every
generated corporate action, and a stale-frame or wrong-side-of-ex read is refused
(Section~\ref{sec:witness}); and each of W1--W4 produces its determined disposition on
replay (Section~\ref{sec:failures}). These restate as minima in
Chapter~\ref{ch:requirements}. The boundary is then closed: what crosses it is recorded,
framed, and replayable, and the conservation the rest of the specification proves rests on
solid ground.
